% Schematic timing of endpoint coordination and the proposed \sys{} path.
% Coordinates are schematic, not measurements. Dashed = control; solid = \data{}.
\begin{figure*}[t]
\centering
\begin{tikzpicture}[x=1cm,y=-1cm,font=\footnotesize]
\path[use as bounding box] (0,0) rectangle (17.5,9.3);
\begin{scope}[shift={(0,4.75)}]
\node[paneltitle] at (.1,.15) {(c) Aligned ranks: endpoint exchanges};
\node at (1.1,.65) {Rank 0}; \node at (5.7,.65) {Rank 1};
\draw[lifeline] (1.1,.9)--(1.1,4.2); \draw[lifeline] (5.7,.9)--(5.7,4.2);
\draw[ctrlarrow] (1.1,1.05)--(5.7,1.65);
\draw[ctrlarrow] (5.7,1.05)--(1.1,1.65);
\node[smalllabel] at (3.4,1.08) {\ready{} all-to-all};
\foreach \y in {1.8,2.04,2.28,2.95} {
\draw[dataarrow] (1.1,\y)--(5.7,{\y+.6});
\draw[dataarrow] (5.7,\y)--(1.1,{\y+.6});}
\draw[ctrlarrow] (1.1,3.62)--(5.7,4.17);
\draw[ctrlarrow] (5.7,3.62)--(1.1,4.17);
\node[smalllabel] at (3.4,3.94) {completion exchange};
\draw[<->] (6.1,1.05)--(6.1,1.65) node[midway,right,align=left] {pre\,$\simeq L$};
\draw[<->] (6.1,3.62)--(6.1,4.17) node[midway,right,align=left] {post\,$\simeq L$};
\end{scope}
\begin{scope}[shift={(8.8,4.75)}]
\node[paneltitle] at (.1,.15) {(d) Aligned ranks: \sys{} overlap};
\node at (.8,.65) {Rank 0}; \node at (3.6,.65) {\sys{}}; \node at (6.4,.65) {Rank 1};
\foreach \x in {.8,3.6,6.4} \draw[lifeline] (\x,.9)--(\x,4.2);
\draw[ctrlarrow] (.8,1.05)--(3.6,1.35); \draw[ctrlarrow] (6.4,1.05)--(3.6,1.35);
\fill[incgreen] (3.6,1.35) circle (.045);
\node[smalllabel,anchor=west] at (3.75,1.08) {all \ready{}};
\foreach \y in {1.2,1.44,1.68,2.35} {
\draw[dataarrow] (.8,\y)--(6.4,{\y+.6});
\draw[dataarrow] (6.4,\y)--(.8,{\y+.6});}
\draw[ctrlarrow] (3.6,2.72)--(6.4,3.02);
\draw[ctrlarrow] (3.6,2.72)--(.8,3.02);
\node[smalllabel,align=center] at (3.6,3.44) {\done{} follows final \data{}\\on each ordered output};
\node[anchor=west] at (.4,4.02) {No sender release wait; no new endpoint exchange.};
\end{scope}
\begin{scope}[shift={(0,0)}]
\node[paneltitle] at (.1,.15) {(a) Unequal ranks: endpoint exchanges};
\node at (1.1,.65) {Rank 0}; \node at (5.7,.65) {Rank 1};
\draw[lifeline] (1.1,.9)--(1.1,4.3); \draw[lifeline] (5.7,.9)--(5.7,4.3);
\draw[ctrlarrow] (1.1,1.0)--(5.7,1.6); \draw[ctrlarrow] (5.7,1.55)--(1.1,2.15);
\node[smalllabel] at (3.4,1.02) {Rank 1 ready later};
\foreach \y in {2.2,2.43,2.66,3.03} \draw[dataarrow] (1.1,\y)--(5.7,{\y+.6});
\foreach \y in {1.65,1.88,2.11,2.48} \draw[dataarrow] (5.7,\y)--(1.1,{\y+.6});
\draw[ctrlarrow] (1.1,3.2)--(5.7,3.8);
\draw[ctrlarrow] (5.7,3.72)--(1.1,4.32);
\node[smalllabel,align=center] at (3.4,4.15) {Readiness and completion\\still traverse between endpoints};
\end{scope}
\begin{scope}[shift={(8.8,0)}]
\node[paneltitle] at (.1,.15) {(b) Unequal ranks: bounded staging};
\node at (.8,.65) {Rank 0}; \node at (3.6,.65) {\sys{}}; \node at (6.4,.65) {Rank 1};
\foreach \x in {.8,3.6,6.4} \draw[lifeline] (\x,.9)--(\x,4.3);
\fill[pattern=north east lines,pattern color=black!35] (3.43,1.43) rectangle (3.77,1.87);
\draw[ctrlarrow] (.8,1.0)--(3.6,1.3);
\draw[ctrlarrow] (6.4,1.55)--(3.6,1.85);
\draw[dataarrow] (.8,1.13)--(3.6,1.43);
\draw[dataarrow] (.8,1.36)--(3.6,1.66);
\draw[dataarrow] (.8,1.59)--(3.6,1.89);
\draw[dataarrow] (3.6,1.9)--(6.4,2.2);
\draw[dataarrow] (3.6,2.13)--(6.4,2.43);
\draw[dataarrow] (3.6,2.36)--(6.4,2.66);
\draw[dataarrow] (.8,2.18)--(3.6,2.48)--(6.4,2.95);
\foreach \y in {1.68,1.91,2.14,2.51} \draw[dataarrow] (6.4,\y)--(.8,{\y+.6});
\draw[ctrlarrow] (3.6,2.72)--(6.4,3.02);
\draw[ctrlarrow] (3.6,2.88)--(.8,3.18);
\node[smalllabel,anchor=east] at (3.2,1.68) {buffer};
\node[smalllabel,align=center] at (3.6,3.6) {Wait for \ready{} at the \sys{} node;\\preserve output service and ordering};
\node[anchor=west] at (.4,4.18) {Skew coverage is not an additional fixed saving.};
\end{scope}
\end{tikzpicture}
\caption{Idealized projection of the two-boundary design, not a measured Direct trace. Solid arrows carry \data{}; dashed arrows carry control. The left-hand reference assumes an exposed tail-signal interval of approximately $L$ after data arrival; actual Direct timing follows Equation~\eqref{eq:tailgain}. \sys{} overlaps readiness with upload and appends ordered \done{} at each destination's tail. Early data is buffered under skew. Required local completion and buffer ownership are retained.}
\Description{Four timing panels compare endpoint and network coordination. The top pair shows unequal readiness and buffering; the lower pair shows aligned readiness as an ideal reference. Solid arrows are data and dashed arrows are control.}
\label{fig:timing}
\end{figure*}
